\section*{Chapter \ref{ch:g5:weather-and-climate} --- Weather and Climate}

\begin{solution}{exo:g5:weather-and-climate:1}
(a) \omterm{def:g1:watching-the-weather:weather}{Weather} (this morning only). (b) \omterm{def:g5:weather-and-climate:climate}{Climate} (every summer). (c) \omterm{def:g1:watching-the-weather:weather}{Weather}
(yesterday afternoon). (d) \omterm{def:g5:weather-and-climate:climate}{Climate} (every month, year after year).
\end{solution}

\begin{solution}{exo:g5:weather-and-climate:2}
The \omterm{def:g5:weather-and-climate:climate}{climate} of a place is the \omterm{def:g1:watching-the-weather:weather}{weather} it usually has, month by month, over
many years. \omterm{def:g1:watching-the-weather:weather}{Weather} scientists average the measurements of 30 years.
\end{solution}

\begin{solution}{exo:g5:weather-and-climate:3}
The coldest month is December (about 6~$^\circ$C; January is almost as
cold). The wettest month is November (about 160~mm). No: the curve never
goes below 0~$^\circ$C.
\end{solution}

\begin{solution}{exo:g5:weather-and-climate:4}
July is the warmest month, at about 5~$^\circ$C. No month is warmer than
10~$^\circ$C.
\end{solution}

\begin{solution}{exo:g5:weather-and-climate:5}
The \omterm{def:g3:continents-and-oceans:equator}{equator}: the hot zone. The \omterm{def:g3:continents-and-oceans:poles}{North Pole}: the cold zone (the northern
one). Halfway between the \omterm{def:g3:continents-and-oceans:equator}{equator} and the \omterm{def:g3:continents-and-oceans:poles}{South Pole}: the southern
temperate zone (it lies between the \omterm{def:g5:weather-and-climate:tropics}{Tropic} of Capricorn and the Antarctic
Circle).
\end{solution}

\begin{solution}{exo:g5:weather-and-climate:6}
Hilo: the hot zone. Phoenix and Seattle: the northern temperate zone.
Utqiaġvik: the northern cold zone, beyond the Arctic Circle.
\end{solution}

\begin{solution}{exo:g5:weather-and-climate:7}
Whole year: $147 + 96 + 106 + 81 + 48 + 37 + 15 + 25 + 41 + 99 + 160 + 145
= 1\,000$ millimetres. December, January and February:
$145 + 147 + 96 = 388$ millimetres, more than a third of the year's \omterm{def:g4:clouds-rain-snow:precipitation}{rain} in
three months.
\end{solution}

\begin{solution}{exo:g5:weather-and-climate:8}
Phoenix: $35 - 14 = 21$ degrees. Hilo: $25 - 22 = 3$ degrees. The seasons
change the temperature much more in Phoenix; in Hilo it is almost the same
all year.
\end{solution}

\begin{solution}{exo:g5:weather-and-climate:9}
A \omterm{def:g5:weather-and-climate:rainforest}{rainforest} \omterm{def:g5:weather-and-climate:climate}{climate}: it is hot all year (27 or 28~$^\circ$C, with no
winter), and a lot of \omterm{def:g4:clouds-rain-snow:precipitation}{rain} falls in every month (at least 160~mm, and more
than 2\,000~mm in the year). The town is most likely in the hot zone, near
the \omterm{def:g3:continents-and-oceans:equator}{equator}.
\end{solution}

\begin{solution}{exo:g5:weather-and-climate:10}
First explanation: the beach is in the hot zone, where it is warm all year,
and the snowy garden in a temperate or cold zone, where January is in
winter. Second explanation: the beach is in the southern half of the Earth,
where January is in summer, and the snowy garden in the northern half,
where January is in winter: the seasons are reversed between the two halves
of the \omterm{def:g3:continents-and-oceans:globe}{globe}.
\end{solution}

\begin{solution}{exo:g5:weather-and-climate:11}
Temperatures: Phoenix is much hotter, about 35~$^\circ$C in July against
about 20~$^\circ$C in August for Seattle, and its January (about
14~$^\circ$C) is already as warm as Seattle's May. \omterm{def:g4:clouds-rain-snow:precipitation}{Rain}: less than 200~mm a year
in Phoenix, a \omterm{def:g5:weather-and-climate:desert}{desert}, against about 1\,000~mm in Seattle. So they do not
have the same \omterm{def:g5:weather-and-climate:climate}{climate}. The zones tell us, roughly, how much the Sun warms
a place, from the \omterm{def:g3:continents-and-oceans:equator}{equator} to the \omterm{def:g3:continents-and-oceans:poles}{poles}; they do not tell us how much \omterm{def:g4:clouds-rain-snow:precipitation}{rain}
falls, nor the effect of the \omterm{def:g2:water-on-earth:ocean}{sea}, of the \omterm{def:g1:landscapes-around-us:landform}{mountains} or of the winds.
\end{solution}
